\begin{afterword}
\summaryhead

The post adds a ``sensitive thingy'' S to the interferometer of ``Configurations and Amplitude.'' S sits on one path and changes from No to Yes when a photon passes. Because every configuration includes S, the two amplitude flows that used to cancel at Detector 1 now go to different configurations, and the photon reaches each detector half the time. The author stresses that this holds whether or not anyone cares about S, looks at it, knows it exists, or could ever learn its state, provided S is moved by more than its original spread. Configurations, the post concludes, are not beliefs; their distinctness is a physical fact.

The second half imagines an early scientist who sees the interference vanish when a sensor is added and concludes that knowledge, or consciousness, decides the result. The author says clues were available, since unread sensors and stray particles do the same, calls it embarrassing that the view survived after the theory predicted that any nudged particle would do, and says it can still be read in textbooks.

\responsehead

The physics of the first half is right. The amplitudes check, the four-way split follows, and the central point, that a physical record of the path removes the interference whether or not anyone reads it, is standard. It holds in every interpretation. The post also states the condition that matters: S must be moved by more than its original spread.

It then drops that condition. The summary that follows says configurations are distinct ``so long as at least one particle in the universe anywhere is in a different position.'' The loss of interference comes in degrees, and a small nudge leaves most of it. The parenthesis states the physics correctly; the summary states it without the condition.

The conclusion, ``Configurations are not belief states,'' does less than it seems to. The experiment refutes the view that what someone knows decides the result. It does not decide whether amplitudes are physical, which is the question interpretations dispute, since collapse theories and epistemic views predict the same outcome. The same goes for the statement that every configuration is about all the particles in the universe. That is an interpretive choice, Everett's, which Bohr rejected.

The history in the second half runs two views together. The imagined scientist holds that what can be known decides the result, and the clues the post lists, an unread sensor and a stray particle, tell against that view. They do not tell against the view the post attaches to von Neumann. Von Neumann did tie collapse to the observer's consciousness, but he also showed that the predictions do not depend on where the line between system and observer is drawn; on his account, too, an unread sensor removes the interference. What he gave to consciousness was the appearance of a single outcome, which these experiments do not test. So the ``embarrassing'' failure to see that any nudged particle would do cannot be charged to von Neumann. The ``Common Wisdom'' the post describes also merges the founders: Bohr, on the \textsc{Stanford Encyclopedia}'s account, never thought of collapse as caused by observers, and treated the interaction with the apparatus as objective. The claim that consciousness views can still be read in textbooks names no textbook. The claim that most theoretical physicists know better fits the one small survey found, in which 6 per cent of foundations researchers gave the observer a special physical role.

In short: correct physics of which-path marking, a summary that drops the condition its own parenthesis gives, and a history whose clues tell against a naive view of knowledge but not against von Neumann's, whom it names.
\end{afterword}
